\documentclass[11pt]{article}

\usepackage[margin=1in]{geometry}
\usepackage{amsmath,amssymb,amsthm,mathtools}
\usepackage[hidelinks]{hyperref}

\newtheorem{theorem}{Theorem}[section]
\newtheorem{lemma}[theorem]{Lemma}
\newtheorem{proposition}[theorem]{Proposition}
\theoremstyle{remark}
\newtheorem{remark}[theorem]{Remark}

\newcommand{\Arg}{\operatorname{Arg}}
\newcommand{\tr}{\operatorname{tr}}

\title{Positive Square Energy for Two $C_5$ Blocks Joined by a Bridge}
\author{}
\date{25 July 2026}

\begin{document}
\maketitle

\begin{abstract}
Let $H$ be the ten-vertex graph formed from two disjoint copies of $C_5$ by
joining one vertex of each copy with a bridge.  We prove that attaching
arbitrary trees at any vertices of this core produces a graph $G$ satisfying
\[
 s^+(G)\geq |V(G)|+s^+(H)-10>|V(G)|.
\]
The proof eliminates the trees by matching belief propagation and compares
the continuous Coulson phase with that of $H$.  The angular comparison is an
exact coefficientwise certificate with $4891$ positive nonconstant terms.
We also factor the characteristic polynomial of $H$ and evaluate $s^+(H)$
exactly.
\end{abstract}

\section{Statement and notation}

All graphs are finite, simple, and undirected.  For adjacency eigenvalues
$\lambda_1,\ldots,\lambda_n$, write
\[
 s^+(G)=\sum_{\lambda_j>0}\lambda_j^2,
 \qquad
 s^-(G)=\sum_{\lambda_j<0}\lambda_j^2.
\]
Let $H$ consist of two vertex-disjoint $5$-cycles and exactly one edge joining
one vertex of the first cycle to one vertex of the second.  Thus $H$ has ten
vertices and eleven edges.

\begin{theorem}\label{thm:main}
Suppose that $G$ is obtained from $H$ by attaching arbitrary rooted trees at
arbitrary core vertices (identifying each tree root with its chosen core
vertex).  If $n=|V(G)|$, then
\begin{equation}\label{eq:main}
 s^+(G)\geq n+s^+(H)-10>n.
\end{equation}
More precisely,
\begin{equation}\label{eq:bare-energy}
 s^+(H)=\frac{31+\sqrt{13}}2-\sqrt5-\rho^2,
\end{equation}
where $\rho$ is the unique root in $(2,\sqrt5)$ of
\begin{equation}\label{eq:rho}
 \rho^3-4\rho-1=0.
\end{equation}
\end{theorem}

The sign in~\eqref{eq:rho} is worth emphasizing.  The characteristic
polynomial has the factor $x^3-4x+1$.  Its negative root is $-\rho$, and
substitution of $x=-\rho$ gives $\rho^3-4\rho-1=0$.  In particular,
$x^3-4x+1$ has no root in $(2,\sqrt5)$; its largest positive root is less
than $2$.  Confusing that positive root with the magnitude of the negative
root changes the sign in~\eqref{eq:rho} and makes the interval assertion
false.

\section{Tree elimination by matching messages}

For a forest $F$, define its signless matching partition
\begin{equation}\label{eq:forest-Z}
 Z_F(t)=\sum_M t^{|V(F)|-2|M|},
\end{equation}
where the sum is over all matchings.  For a graph with positive vertex
activities $a=(a_v)$, use the weighted version
\begin{equation}\label{eq:weighted-Z}
 Z_F(a)=\sum_M\prod_{v\text{ unmatched by }M}a_v.
\end{equation}
All edge activities are one.

Orient every branch outside the core toward the core.  If $u\to v$ is an
oriented branch edge, let $T_{u\to v}$ be the component below $u$, including
$u$, and put
\begin{equation}\label{eq:message-definition}
 m_{u\to v}(t)=
 \frac{Z_{T_{u\to v}}(t)}{Z_{T_{u\to v}-u}(t)}.
\end{equation}

\begin{lemma}[Matching belief propagation]\label{lem:BP}
The messages satisfy
\begin{equation}\label{eq:BP}
 m_{u\to v}(t)=t+\sum_{w\text{ child of }u}
                    \frac1{m_{w\to u}(t)}\geq t>0.
\end{equation}
After all branches are eliminated, each core vertex has activity
\begin{equation}\label{eq:activity}
 a_v(t)=t+\sum_{u\text{ attached at }v}\frac1{m_{u\to v}(t)}
       =t+y_v(t),
 \qquad y_v(t)\geq0.
\end{equation}
Every term in the grouped Sachs expansion has the same positive real factor
\begin{equation}\label{eq:K}
 K(t)=\prod_{u\text{ adjacent to the core}}Z_{T_{u\to v}}(t)>0.
\end{equation}
This remains true for a Sachs term in which either entire pentagon is
deleted.
\end{lemma}

\begin{proof}
Split a matching below $u$ according to whether $u$ is unmatched or is
matched to one of its children, and factor the child-subtree partitions.
This proves~\eqref{eq:BP}; induction from a leaf gives the inequality.
For a branch rooted next to a core vertex, extracting its full partition
changes the weight of using the edge into the branch to
$Z_{T-u}/Z_T=1/m_{u\to v}$.  The mutually exclusive choices over incident
branches give~\eqref{eq:activity}, while independence gives~\eqref{eq:K}.

If a core vertex is deleted as part of a Sachs cycle, every branch attached
there becomes a disconnected full tree and still contributes its factor
$Z_T$ to $K(t)$.  Retained core vertices have the same effective activities
as before.  Hence deleting a whole pentagon does not leave an additional
ratio of branch partitions.
\end{proof}

\section{The weighted bridge core}

Let
\[
 \phi_G(z)=\det(zI-A(G)),
 \qquad
 \Psi_G(t)=i^{-n}\phi_G(it)=\prod_{j=1}^n(t+i\lambda_j),
 \quad t>0.
\]
In the grouped Sachs expansion, a $5$-cycle has multiplier
$-2i^{-5}=2i$.  The two cycles are disjoint, so their joint selection has
multiplier $(2i)^2=-4$.

Root one weighted pentagon at its bridge endpoint and list its other
activities consecutively as $x_1,x_2,x_3,x_4$.  Define
\begin{align}
 P(x_1,x_2,x_3,x_4)
 &=x_1x_2x_3x_4+x_3x_4+x_1x_4+x_1x_2+1,\label{eq:P}\\
 B(x_1,x_2,x_3,x_4)
 &=x_2x_3x_4+x_2+x_4+x_1x_2x_3+x_1+x_3.\label{eq:B}
\end{align}
Here $P$ is the matching partition after deleting the root.  If the root
activity is $a_0$, the full pentagon partition is
\begin{equation}\label{eq:Q}
 Q=a_0P+B.
\end{equation}
For the left pentagon denote the deleted-root and full partitions by $p,q$;
for the right one denote them by $s,r$, respectively.

\begin{proposition}[Exact bridge formula]\label{prop:bridge}
With the activities~\eqref{eq:activity},
\begin{equation}\label{eq:bridge-formula}
 \boxed{\quad
 \frac{\Psi_G(t)}{K(t)}=(qr+ps-4)+2i(q+r).
 \quad}
\end{equation}
\end{proposition}

\begin{proof}
Splitting matchings according to whether the bridge is unused or used gives
$qr+ps$: in the second case both bridge endpoints are deleted.  Selecting
the left Sachs cycle leaves the full right pentagon and contributes $2ir$;
selecting the right one contributes $2iq$.  Selecting both contributes
$-4$.  Lemma~\ref{lem:BP} supplies the common factor $K(t)$.
\end{proof}

For bare activities, set
\begin{equation}\label{eq:bare-PQ}
 P_0=t^4+3t^2+1,
 \qquad Q_0=t^5+5t^3+5t.
\end{equation}
Introduce the real and imaginary coordinates
\begin{align*}
 X&=qr+ps-4,&Y&=2(q+r),\\
 X_0&=Q_0^2+P_0^2-4,&Y_0&=4Q_0.
\end{align*}
Both $Y$ and $Y_0$ are positive for $t>0$.  We therefore use the
upper-half-plane branch
\begin{equation}\label{eq:phases}
 \Theta_G(t)=\operatorname{atan2}(Y,X),
 \qquad
 \Theta_H(t)=\operatorname{atan2}(Y_0,X_0),
 \qquad 0<\Theta_G,\Theta_H<\pi.
\end{equation}
In particular,
\begin{equation}\label{eq:bare-phase}
 \Theta_H(t)=\operatorname{atan2}
 \bigl(4Q_0,Q_0^2+P_0^2-4\bigr).
\end{equation}
An ordinary one-argument arctangent cannot replace this expression globally:
\[
 X_0=t^{10}+11t^8+41t^6+61t^4+31t^2-3
\]
is negative for all sufficiently small positive $t$.

\section{Coefficientwise angular certificate}

\begin{proposition}\label{prop:certificate}
For all activities $a_v=t+y_v$ with $t>0$ and $y_v\geq0$,
\begin{equation}\label{eq:cross-product}
 Y_0X-YX_0\geq0.
\end{equation}
Consequently,
\begin{equation}\label{eq:phase-comparison}
 0<\Theta_G(t)\leq\Theta_H(t)<\pi.
\end{equation}
The cross-product inequality is an equality at the bare specialization
$y_0=\cdots=y_9=0$.
\end{proposition}

\begin{proof}
After canceling a positive factor $2$, inequality~\eqref{eq:cross-product}
is
\begin{equation}\label{eq:certificate-polynomial}
 2Q_0(qr+ps-4)-(q+r)(Q_0^2+P_0^2-4)\geq0.
\end{equation}
Substitute $a_v=t+y_v$ at the ten core vertices, use
\eqref{eq:P}--\eqref{eq:Q}, and expand over $\mathbb Z$ in the ordered
variables $(t,y_0,\ldots,y_9)$.  The exact expansion has $4891$ nonzero
monomials.  Every coefficient is positive, the minimum is $1$, the maximum is
$289$, and every monomial contains at least one $y_v$.  Thus the polynomial
is nonnegative on the nonnegative orthant and has zero bare specialization.
The exhaustive symbolic verification and its canonical digest are recorded
in Appendix~\ref{app:certificate}.

It remains to translate the cross product without losing the quadrant.  From
\eqref{eq:phases}, the difference
$\Theta_G-\Theta_H$ belongs to $(-\pi,\pi)$, and
\[
 \sin(\Theta_G-\Theta_H)
 =\frac{YX_0-XY_0}{\sqrt{X^2+Y^2}\sqrt{X_0^2+Y_0^2}}
 \leq0.
\]
On $(-\pi,\pi)$ this implies $\Theta_G-\Theta_H\leq0$.  No sign assumption
on $X$ or $X_0$ and no ordinary-arctangent division is used.
\end{proof}

\section{Phase integration and the bare spectrum}

For the continuous argument in~\eqref{eq:phases}, the signed square-energy
Coulson identity is
\begin{equation}\label{eq:Coulson}
 s^+(F)-s^-(F)
 =-\frac4\pi\int_0^\infty t\Theta_F(t)\,dt.
\end{equation}
It follows by applying the scalar identity to the factors
$t+i\lambda_j$; the trace relation $\sum_j\lambda_j=0$ cancels the
nonintegrable first-order term at infinity.  Proposition~\ref{prop:certificate}
and~\eqref{eq:Coulson} give
\begin{equation}\label{eq:difference-comparison}
 s^+(G)-s^-(G)\geq s^+(H)-s^-(H).
\end{equation}

We now evaluate the right side exactly.  A direct determinant calculation
gives
\begin{equation}\label{eq:charpoly}
 \phi_H(x)=(x-1)(x^2-x-3)(x^2+x-1)^2(x^3-4x+1).
\end{equation}
The cubic has two positive roots $\alpha,\beta$ and one negative root
$-\rho$.  Indeed, $f(x)=x^3-4x+1$ has one root in each of
$(-3,-2)$, $(0,1)$, and $(1,2)$.  Vieta's relations give
\[
 \alpha+\beta=\rho,
 \qquad
 \alpha\beta=\rho^2-4,
 \qquad
 \alpha^2+\beta^2=8-\rho^2.
\]
The other positive eigenvalues read from~\eqref{eq:charpoly} are
\[
 1,
 \qquad \frac{1+\sqrt{13}}2,
 \qquad \frac{-1+\sqrt5}2\quad\text{with multiplicity two}.
\]
Squaring and summing proves~\eqref{eq:bare-energy}.

For completeness, the definition of $\rho$ and the strict slack can be
verified without decimals.  Let $g(x)=x^3-4x-1$.  Then
\[
 g(2)=-1<0,
 \qquad
 g(\sqrt5)=\sqrt5-1>0,
 \qquad
 g'(x)=3x^2-4>0\quad(x>2).
\]
Hence $g$ has exactly one root $\rho$ in $(2,\sqrt5)$, and $\rho^2<5$.
Therefore
\begin{align}
 s^+(H)
 &>\frac{21+\sqrt{13}}2-\sqrt5>10.\label{eq:slack}
\end{align}
The last inequality is exact: it is equivalent to
$(1+\sqrt{13})/2>\sqrt5$, and both sides are positive; after moving $1$ and
squaring it reduces to $13>(2\sqrt5-1)^2=21-4\sqrt5$, equivalently
$\sqrt5>2$.

Finally, $G$ is connected and bicyclic, so it has $n+1$ edges.  Thus
\begin{equation}\label{eq:trace}
 s^+(G)+s^-(G)=\tr A(G)^2=2n+2,
 \qquad
 s^+(H)+s^-(H)=22.
\end{equation}
Averaging~\eqref{eq:difference-comparison} with~\eqref{eq:trace} yields
\[
 s^+(G)\geq n+1+\frac{s^+(H)-s^-(H)}2
          =n+s^+(H)-10.
\]
The strict inequality in~\eqref{eq:main} follows from~\eqref{eq:slack}.
Equivalently, the exact bare phase integral is
\[
 \int_0^\infty t\Theta_H(t)\,dt
 =\frac\pi2\bigl(11-s^+(H)\bigr)
 =\frac\pi2\left(
 \rho^2+\sqrt5-\frac{9+\sqrt{13}}2\right)<\frac\pi2.
\]

\appendix
\section{Exact certificate and reproducibility}\label{app:certificate}

The certificate script is
\begin{center}
\texttt{positive-square-energy/experiments/\allowbreak
c5\_bridge\_phase\_certificate.py}.
\end{center}
Run it from the repository root with
\begin{verbatim}
python positive-square-energy/experiments/c5_bridge_phase_certificate.py
\end{verbatim}
It constructs~\eqref{eq:certificate-polynomial} over $\mathbb Z$, checks the
bare characteristic-polynomial factorization~\eqref{eq:charpoly}, and asserts
all of the following for the coefficient expansion:
\begin{enumerate}
\item there are exactly $4891$ nonzero monomials;
\item every nonzero coefficient is positive, with minimum $1$ and maximum
      $289$;
\item every monomial contains at least one of $y_0,\ldots,y_9$;
\item the bare specialization is zero; and
\item the SHA-256 digest of the canonical ordered stream
      \texttt{monomial:coefficient}, one term per line, is
\end{enumerate}
\begin{center}
\texttt{365e18dedf9032fbcdb88af83d033f0651f02412c796b4f1dfde04152a478af1}.
\end{center}
The digest commits to the complete coefficient list rather than to a
numerical sampling of activities.

\end{document}
